% !TeX root = ../初等数学笔记.tex
% 由主文档通过 \input 载入；请编译主文档。
\setcounter{chapter}{14} % 保留原书章号。
\chapter{几何变换、轨迹与作图}
\label{chap:em-transformations}
几何变换通过移动或缩放整个图形，把难以直接比较的量放到同一位置。轨迹和作图则从相反方向出发：给定条件，寻找全部满足条件的点或构造出所需对象。

\section{平移、旋转与轴对称}
\emph{平移}（translation）把所有点沿同一方向移动同一距离；\emph{旋转}（rotation）固定中心 $O$，把各射线 $OP$ 转过同一有向角，并保持 $OP$ 长度；\emph{轴对称}（reflection）把点映到关于固定直线的镜像位置，连接点与像的线段被该直线垂直平分，轴上点保持不动。

这些变换都保持任意两点之间的距离，因而保持线段长度、角的大小、全等关系与面积。半周旋转就是中心对称。轴对称会颠倒平面定向，平移和旋转保持定向；因此“保持距离”不等于“保持顺逆时针次序”。变换复合依次执行，顺序一般不能交换，例如先旋转再沿固定方向平移通常不同于反过来操作。

\section{位似与相似变换}
给定中心 $O$ 和实数 $k\ne0$，\emph{位似}（homothety）把 $P$ 映为 $P'$：当 $k>0$ 时两点在同一射线 $OP$ 上且 $OP'=kOP$；当 $k<0$ 时在相反射线上且 $OP'=|k|OP$，中心仍映为自身。由相似三角形可得 $P'Q'=|k|PQ$，不经过中心的直线的像与原线平行。

正位似是缩放，负位似是缩放与中心对称的复合。与保持距离的变换复合后得到一般相似变换，它保持角的大小，把所有长度统一乘以同一正因子。圆的像仍为圆，半径乘以 $|k|$；面积乘以 $k^2$。比例 $k=0$ 会把所有点压到一点，不属于本节非退化相似变换。

\section{轨迹的必要性与充分性}
\emph{轨迹}（locus）是所有满足指定条件的点组成的集合。候选图形必须经过双向检验：每个合法点都在候选图形上，候选图形的每个允许点也满足原条件。只做到第一步，可能把真正轨迹扩大。

到定点 $O$ 距离为 $r>0$ 的点集是圆；到两个不同定点 $A,B$ 等距的点集是线段 $AB$ 的垂直平分线。到两条相交直线距离相等的点集是两条夹角平分线的并，而不是仅一条内角平分射线；若点被限制在某个角内部，才取其中对应分支。两不同平行线的等距点集则是它们之间的平行中线。

\section{尺规作图与存在条件}
尺规作图只允许使用无刻度直尺作两已知点连线，以及用圆规以已知点为圆心、已知长度为半径作圆，再取已作对象的交点。证明包括分析、构造、验证和解数讨论；能在图上估计一个点，不等于已给出尺规构造。

例如作 $AB$ 的垂直平分线，以 $A,B$ 为圆心取同一半径 $r>AB/2$ 作圆，交于 $P,Q$，连 $PQ$。两交点到 $A,B$ 等距，故均在垂直平分线上；两点确定直线，所作即目标。要求 $r>AB/2$ 是为了得到两个不同交点。

给定边长 $a,b,c>0$ 作三角形，可先作 $AB=c$，再取以 $A,B$ 为圆心、半径分别为 $b,a$ 的两圆交点。非退化交点存在当且仅当 $|a-b|<c<a+b$；两交点关于 $AB$ 对称，得到两种放置但同一个全等类型。等号情形共线，严格不等式失败则不可作成所需三角形。

\section{辅助构造、最短路径与方法比较}
\begin{theorem}[反射求最短折线]
设 $A,B$ 位于直线 $\ell$ 同一侧且都不在 $\ell$ 上。令 $A'$ 为 $A$ 关于 $\ell$ 的对称点，则 $P$ 遍历 $\ell$ 时，$AP+PB$ 的最小值为 $A'B$，在 $\ell$ 与线段 $A'B$ 的交点处取得。
\end{theorem}
\begin{proof}
对每个 $P\in\ell$，对称性给出 $AP=A'P$，所以 $AP+PB=A'P+PB\ge A'B$。$A',B$ 在直线两侧，连接线段恰穿过直线一次；在该交点处两段连接成一条直线且次序正确，三角不等式取等。
\end{proof}
若 $P$ 被进一步限制在直线的某个区段，交点可能不合法，须重新研究边界；不能只把无限直线上的答案照搬。辅助线的作用也是把原图形嵌入更容易使用全等、相似或三角不等式的结构，而不是无目的地增加线条。

\section{例题、习题与学习检查}
\begin{example}
在直角坐标图中，$A=(0,1),B=(4,3)$，折线 $APB$ 的转折点 $P$ 在横轴上。求最短总长。
\end{example}
\begin{solve}
把 $A$ 反射为 $A'=(0,-1)$，由勾股定理 $A'B=\sqrt{4^2+4^2}=4\sqrt2$。连接线的竖直变化从 $-1$ 到 $3$，与横轴相交时完成全程的四分之一，故 $P=(1,0)$，最短值确实可达。
\end{solve}
\begin{enumerate}
\item 位似比为 $-3$，长度和面积分别乘以多少？答：$3$ 与 $9$。
\item 给定边长 $2,3,6$ 能否作三角形？答：不能，因为 $2+3<6$。
\item 到两个不同点等距且到其中一个点距离为 $r$ 的点最多有几个？答：两个，由垂直平分线与圆的交点决定，也可能为零或一个。
\end{enumerate}

